\chapter{Curvatura, Gauss--Bonnet y superficies}
\label{cap:curvatura-superficies}

La geometría hiperbólica no se limita a los dibujos del disco y del semiplano. También aparece en superficies cuya curvatura es negativa y en cocientes del plano por grupos de isometrías. El teorema de Gauss--Bonnet relaciona curvatura local, área y topología \cite{ratcliffe2006,stillwell1992}.

\section{Curvatura constante}

En una superficie riemanniana, la curvatura gaussiana mide cómo se aparta la geometría infinitesimal de la euclidiana. En el plano hiperbólico normalizado vale \(K=-1\). Si se escala la métrica por un factor constante, cambia el valor numérico de la curvatura, aunque la geometría resultante sigue siendo de curvatura negativa constante.

\section{Gauss--Bonnet}

Para una región compacta \(R\) con borde geodésico a trozos, Gauss--Bonnet relaciona la curvatura interior y los ángulos exteriores:
\begin{equation}
  \int_R K\,dA+\sum_j(\pi-\alpha_j)=2\pi\chi(R).
  \label{eq:gauss-bonnet-borde}
\end{equation}
Aquí \(\chi(R)\) es la característica de Euler y \(\alpha_j\) son los ángulos interiores en las esquinas. Para un disco topológico, \(\chi(R)=1\); si además \(K=-1\), la fórmula se reduce a la ley de área usada en el \cref{cap:triangulos-hiperbolicos}.

\begin{theorem}[Área de una superficie hiperbólica cerrada]
Si \(S\) es una superficie cerrada, orientable y de curvatura constante \(-1\), entonces
\begin{equation}
  \operatorname{Area}(S)=-2\pi\chi(S)=4\pi(g-1),
  \label{eq:area-superficie-cerrada}
\end{equation}
donde \(g\geq 2\) es el género de \(S\).
\end{theorem}

\begin{proof}
Gauss--Bonnet para una superficie cerrada afirma que \(\int_S K\,dA=2\pi\chi(S)\). Como \(K=-1\), el lado izquierdo es \(-\operatorname{Area}(S)\), así que el área es \(-2\pi\chi(S)\). Para una superficie orientable de género \(g\), \(\chi(S)=2-2g\), lo que da la segunda igualdad.
\end{proof}

\section{Construir superficies por identificación}

Una manera de obtener una superficie cerrada es identificar lados de un polígono hiperbólico según un patrón prescrito. Los lados identificados heredan la misma longitud, y las condiciones angulares garantizan que alrededor de cada vértice la superficie no tenga una singularidad cónica. Para una explicación más completa de los cocientes y variedades, véase \cite{ratcliffe2006}.

\begin{example}[Género dos]
Una superficie cerrada de género \(2\) tiene \(\chi=-2\), de modo que \cref{eq:area-superficie-cerrada} da área \(4\pi\). El área depende del género y de la normalización de curvatura, no de la forma concreta del polígono usado para construirla.
\end{example}
